\section{Algorithm}\label{sec:algorithm}

\paragraph{Purification.}
Given a QF\_NIA formula $\varphi$, the first step is \emph{purification}: every
nonlinear subterm $t$ (a product whose factors are not all constants, or an
integer division/modulo with a non-constant divisor) is replaced by a fresh
integer constant $\pure{t}$, called a \emph{pure}.
Equal subterms share the same pure.
The result $\hat\varphi$ is a QF\_LIA formula over the original variables and
the new pure constants.

\paragraph{Main loop.}
\Cref{alg:solve} shows the top-level procedure.
We maintain a set $\mathcal{A}$ of linear \emph{axioms}---constraints that
relate each pure $\pure{t}$ to its factors---accumulated across iterations.
Each iteration calls a QF\_LIA solver on $\hat\varphi \land \bigwedge\mathcal{A}$.
If the LIA problem is unsatisfiable, so is the original QF\_NIA formula.
Otherwise we obtain a model $\mu$ and check whether it is consistent with the
nonlinear semantics (\Cref{alg:checknia}).
If yes, $\mu$ witnesses satisfiability.
If no, new axioms are added and the loop continues.

\begin{algorithm}[t]
\caption{QF-NIA-SOLVE($\varphi$)}\label{alg:solve}
\DontPrintSemicolon
$\hat\varphi,\, P \leftarrow \textrm{PURIFY}(\varphi)$\;
$\mathcal{A} \leftarrow \emptyset$\;
\Repeat{}{
  $\mathit{res} \leftarrow \textrm{CHECK-LIA}(\hat\varphi \land \bigwedge\mathcal{A})$\;
  \lIf{$\mathit{res} = \mathit{unsat}$}{\Return $\mathit{unsat}$}
  $\mu \leftarrow \textrm{model}(\mathit{res})$\;
  \lIf{\upshape CHECK-NIA$(P,\mu,\mathcal{A})$}{\Return $\mathit{sat},\, \mu$}
}
\end{algorithm}

\begin{algorithm}[t]
\caption{CHECK-NIA($P, \mu, \mathcal{A}$)}\label{alg:checknia}
\DontPrintSemicolon
\lIf{\upshape QUICK-CHECK$(P,\mu)$}{\Return \upshape true}
$I \leftarrow \textrm{IMPLICANT}(\hat\varphi, \mu)$\tcp*{literals of $\hat\varphi$ satisfied by $\mu$}
$R \leftarrow \{\,p \in P \mid p \in I,\; \mu(p) \neq \mu(\mathrm{term}(p))\,\}$\tcp*{failing pures}
\lIf{$R = \emptyset$}{\Return \upshape true}
add lazy congruence axioms for $R$ to $\mathcal{A}$\;
\ForEach{$p \in R$}{
  $\mathcal{A} \leftarrow \mathcal{A} \cup \textrm{AXIOMS}(p, \mu)$\tcp*{all violated axioms for $\mathrm{term}(p)$}
}
\Return \upshape false\;
\end{algorithm}

\paragraph{Quick check.}
Before collecting failing pures, we perform a fast three-valued evaluation of
$\hat\varphi$ under $\mu$, substituting each pure $\pure{t}$ with the actual
value $\mu(t)$ (evaluated under NIA semantics).
If the result is \TRUE, no axioms are needed and we return immediately.

\paragraph{Implicant-based targeting.}
Rather than scanning all pures in $\hat\varphi$, we first compute an
\emph{implicant}: a minimal subset $I$ of literals of $\hat\varphi$ that are
jointly sufficient for $\mu \models \hat\varphi$.
We then restrict attention to pures appearing in literals of $I$ that are
falsified under NIA semantics, ignoring pures in satisfied literals.
This reduces the number of axioms generated per iteration and focuses
refinement on the parts of the formula that actually need it.

\paragraph{Axiom generation.}
For each failing pure $\pure{t}$, we add to $\mathcal{A}$ every axiom in our
axiom set for $t$ that is violated by $\mu$.
Axiom classes are described in \Cref{sec:axioms}; they cover sign and zero
conditions, secant-based linear bounds for monomials $x^k$, bounds for mixed
products $x^k y^l$, tangent-plane lemmas, and integer division and modulo
constraints.
All violated axioms are added in a single pass---there is no round structure.
Lazy congruence axioms (requiring $\pure{t_1} = \pure{t_2}$ whenever
$t_1 = t_2$ holds in $\mu$ but $\mu(\pure{t_1}) \neq \mu(\pure{t_2})$) are
handled separately before per-pure axiom generation.

\paragraph{Comparison with~\cite{cimatti2018sat}.}
The overall CEGAR skeleton is the same.
The key differences are as follows.

\emph{Abstract domain.}
Cimatti et al.\ replace each product $x \cdot y$ with an uninterpreted
function symbol $f(x,y)$, yielding a QF\_UFLIA abstraction; congruence
($f(a,b) = f(a,b)$ for equal arguments) is then automatic.
We use plain fresh constants instead, so the abstract domain is QF\_LIA, and
congruence must be added as explicit axioms when violated.

\emph{Axiom generation.}
Cimatti et al.\ generate lemmas in three sequential rounds (basic axioms,
proportionality, tangent-plane), stopping as soon as a round produces at
least one lemma.
We add all violated axioms in a single pass, regardless of type.

\emph{Failing-term selection.}
Cimatti et al.\ scan all constraints of $\varphi$ falsified by $\mu$ and
collect every product $x \cdot y$ in those constraints.
We use an implicant of $\hat\varphi$ under $\mu$ and restrict to pures in
literals that fail under NIA semantics, which is more targeted.

\emph{Axiom set.}
The axioms in~\cite{cimatti2018sat} cover sign, zero, neutrality,
proportionality, and tangent-plane conditions for binary products.
We retain sign and tangent-plane axioms and add secant-based bounds for
monomials $x^k$ and mixed products $x^k y^l$, described in \Cref{sec:axioms}.
For tangent-plane lemmas we optionally employ the \emph{frontier} strategy
of~\cite{cimatti2017tacas}: a per-pure bounding box tracks all model values
seen so far, and whenever a new point extends the box, extra tangent planes
are instantiated at the two mixed corners to ensure both upper and lower
bounds exist in every region.
